\section{The inverse-polynomial option and completion of the proof}
\label{sec:rare-overlap}

In the remaining alternative of \Cref{prop:status-preparation},
compatible pairs of leaves have mass \(\Omega(N^{-4})\).
We will prove an expected overlap lower bound of the same order.
The compatible event can have zero mass in the common limit of
\Cref{sec:common-limit}, so we need exceptional-set estimates smaller
than \(N^{-4}\).

The endpoint identities in this alternative are exact:
\begin{equation}
 p_i=T(\lambda_i,\cdot)+a.
\label{eq:rare-exact-prediction}
\end{equation}
The bit \(c=a(\lambda_1)+a(\lambda_2)\) is common to the prepared law,
and the old marked tensors and their evaluations on the fixed cover
have been pinned.  Compatibility is therefore determined by the two
leaves before their orientations are sampled.  We shall use the exact
identities to choose one generic atom per cross edge, with no additional
\(p_i\)-status filter.

First we prove a hit estimate for each prescribed large reference set,
with a uniform superpolynomial bound on the exceptional orientation
mass.  This estimate uses the mixed-law bounds.  The original prepared
leaves' histogram bounds then supply a finite family of large sets on
each B-leaf.  Choosing that family before the independent A-leaf is
sampled lets us apply the hit estimate to all its members.

\subsection{Generic base queries}

In this Section a \emph{base query} is a successful-query format with
at most four atom positions in one unit, with labels distinct at each
endpoint.  Every position uses the full generic flavor, a prescribed
possible diagonal bit \(q\), a feasible tester summary, and a prescribed
assignment of the basis bits
\(T(x,\cdot)|_{C_i}\).  It also imposes the reference off-position
Gram zeros and a consistent numerical specification of the binary
evaluations used in \Cref{thm:gradient-realization}.
There is no additional \(p_i\)-status restriction and, for now, no
orientation-only filter.

The point inputs are independent, each normalized conditional on its
prescribed \(q\).  For a unit orientation \(\omega\), write
\(Q_{\omega,\mathfrak r}\) for the unnormalized measure of its accepted
key tuples in format \(\mathfrak r\).  Its total mass is at most one.
The reference space for its shape is
\((\cK_{\mathfrak r,n},\nu_{\mathfrak r,n})\).
The request is \emph{feasible} at \(\omega\) when its coordinate format,
labels, summaries and assignments are allowed there.  A basis of
dimension smaller than \(K^2\) may be padded by zero coordinates;
feasibility then requires the corresponding requested bits to be zero.
There are only finitely many numerical formats and requests, although
their effective tensor bases can depend on the orientation and its leaf.

Let
\begin{equation}
 \beta=2^{-K^2-1}4^{-(g^2+1)}.
 \label{eq:rare-base-beta}
\end{equation}
This is the uniform generic base-status fraction from
\Cref{lem:unary-positivity}; see \eqref{eq:uniform-base-fraction}.
The constant applies to every feasible generic tester summary and
basis assignment on each fixed common key cell of positive mass,
independently of that cell's mass.

Let \(\kappa_{\mathfrak r}>0\) be the product-mixing constant of
\Cref{lem:binary-mixing} for a consistent format, and take
\(\kappa_{\min}>0\) to be the minimum over the finite collection in use.
If necessary include the finitely many coordinate formats separately.
Put \(s_{\max}=4\).

\begin{lemma}[Uniform large-set hit estimate]
\label{lem:rare-large-set-hit}
Fix \(\delta>0\), and consider a sequence of unit laws whose mixed
joint densities are at most \(2^{O(N)}\) relative to \(\mu^2\), whose
two marginals are at most \(M2^{o(N)}\mu\), and whose marked effective
spaces have the rank and dimension bounds of \Cref{lem:unary-positivity}.
The implied joint-density constant is uniform.
Set
\begin{equation}
 c_{\mathrm{hit}}(\delta)
   =\frac14\,\kappa_{\min}\beta^{s_{\max}}\delta.
 \label{eq:rare-hit-threshold}
\end{equation}
For every fixed \(p>0\),
\begin{equation}
 \sup_{\mathfrak r,W_n}
 \PP\!\left(
 \begin{array}{l}
 \mathfrak r\text{ is feasible at }\omega,\\
 Q_{\omega,\mathfrak r}(W_n)<c_{\mathrm{hit}}(\delta)
 \end{array}\right)
 =o(n^{-p}),
 \label{eq:rare-superpolynomial-tail}
\end{equation}
where the supremum is over the finitely many numerical requests
and deterministic measurable sets
\(W_n\subseteq\cK_{\mathfrak r,n}\) satisfying
\(\nu_{\mathfrak r,n}(W_n)\ge\delta\).
Each choice of \(W_n\) has its own exceptional orientations.
\end{lemma}

\begin{proof}
Suppose the conclusion fails for some \(p\).  Along a subsequence
there are deterministic sets \(W_n\), requests \(\mathfrak r\), and
a fixed \(c>0\) for which the displayed bad event has probability
at least \(c n^{-p}\).  By a finite subdivision, fix the numerical
format and request along this subsequence.  Restrict the unit law
to this bad event and normalize it.
The additional density factor is at most \(c^{-1}n^p\).
Since \(N=M_0n\) with \(M_0\) fixed, its logarithm is
\(O(\log n)=o(n)=o(N)\).  The restricted mixed law therefore
still satisfies the hypotheses of \Cref{lem:typicality,thm:weak-product},
and the uniform profile estimates in \Cref{lem:unary-positivity}.

We give the measure-identification argument for this restricted law.
It uses no density bound on its individually normalized old leaves.
On each single-key space choose countably many common bit tests,
and on the tuple-key space choose countably many bit tests containing
\(\1_{W_n}\).  Include the unary factors in weak product approximations
to every tuple cylinder at accuracies tending to zero, their unary
pullbacks, and the resulting finite Boolean combinations.
This countable collection is obtained by successive closure.
At every fixed stage its complexity is bounded independently of \(n\).

Pass to a subsequence on which the reference signature laws,
the bounded coefficients in these approximations, and the laws of
all empirical unary status arrays converge.
The arrays have coordinates in \([0,1]\), indexed by finite cylinders
and the finite status choices; hence this is ordinary compactness
on a countable product of compact spaces.
Write \(\Xi_j\) for the single-key signature space at position \(j\),
and \(\rho\) for the tuple-test reference law.
By \Cref{lem:reference-product}, the unary projections of \(\rho\)
are independent with the respective laws on \(\Xi_j\).
The limiting \(W\)-coordinate is a bit, and
\(\rho(W)\ge\delta\).

Let \(\theta\) denote a limiting orientation status array.
Its unfiltered unary measure equals the single-key reference measure
on every cylinder, by \Cref{lem:typicality}.
The requested unary base-status measure therefore has a density
\(h_j^\theta\) relative to that reference measure.
Versions can be chosen measurably by conditional averages on the
increasing finite cylinder partitions.
The simultaneous positive-fraction estimates on all fixed
positive-reference-mass cylinders give
\begin{equation}
             \beta\le h_j^\theta\le1
 \quad\text{almost everywhere, for almost every }\theta.
 \label{eq:rare-unary-density-lower}
\end{equation}
Indeed these inequalities hold on the generating cylinder algebra
and hence on its generated sigma-algebra by measure domination.
All requests are feasible in the restricted law, so their limiting
coordinate formats and summary assignments are feasible as well.

Also pass to a limit of the unnormalized accepted tuple measures on
the tuple-test cylinders.  Denote it by \(\tau\).
We claim that
\begin{equation}
 \frac{d\tau}{d\rho}
   =\kappa_{\mathfrak r}\,
      \E_\theta\prod_{j=1}^{s}h_j^\theta(\xi_j),
 \qquad s\le s_{\max}.
 \label{eq:rare-one-law-density}
\end{equation}
Here \(\xi_j\) is the unary projection of the tuple signature.
This is the one-law version of the density identification in
\Cref{lem:limit-density}; the details needed here are as follows.

Fix a tuple cylinder and an accuracy \(\epsilon>0\).
By \Cref{thm:weak-product}, its prelimit indicator can be replaced
on actual successful queries by a bounded finite sum of products
of global unary key functions, with error at most \(\epsilon\)
outside a vanishing set of orientations.
This comparison allows the requested unary status indicators as
orientation-dependent unary weights.
For each structured product, \Cref{lem:binary-mixing} gives its
successful integral as \(\kappa_{\mathfrak r}\) times the product
of its separate unary status integrals, with vanishing error.
Averaging over the restricted law preserves these errors:
the relevant estimates use the mixed density and marginal bounds
already checked.

The unary integrals and finite numerical metadata are coordinates of
the status arrays.  Their finite products and averages therefore pass
to the limit.  Independence of the unary projections of \(\rho\)
identifies the result with the structured test against the right side
of \eqref{eq:rare-one-law-density}.
The reference part of \Cref{thm:weak-product} bounds the error of the
same replacement against arbitrary products of unary functions
bounded by one.  Passing this assertion first for cylinder-simple
functions and then by \(L^1\) approximation permits the functions
\(h_j^\theta\) themselves.
After averaging \(\theta\), the discrepancy on the chosen cylinder
is bounded by a constant multiple of \(\epsilon\).
Send \(\epsilon\) to zero.  A countable generating algebra of
cylinders proves \eqref{eq:rare-one-law-density}.

In particular, the limiting mean of
\(Q_{\omega,\mathfrak r}(W_n)\) is
\[
 \tau(W)
 =\kappa_{\mathfrak r}
   \int_W \E_\theta\prod_{j=1}^{s}h_j^\theta(\xi_j)\,d\rho
 \ge\kappa_{\min}\beta^{s_{\max}}\delta.
\]
But each restricted orientation had hit mass smaller than
\(c_{\mathrm{hit}}(\delta)\), contradicting
\eqref{eq:rare-hit-threshold}.
This proves \eqref{eq:rare-superpolynomial-tail}.
\end{proof}

\begin{remark}
\label{rem:rare-restriction-caps}
A restriction of polynomial total mass need not preserve a uniform
density cap for every separately normalized old leaf.
\Cref{lem:rare-large-set-hit} does not require that conclusion:
its proof uses mixed-law typicality and weak product testing, not
the per-leaf net lemma or the two-leaf covariance projection.
The leaf estimates used next concern the original prepared leaves.
Also, the proof of \Cref{lem:unary-positivity} may choose a larger
fixed number of transformations for a smaller positive key cell.
This changes an asymptotic onset and a constant prefactor, not
\(\beta\), \(J\), or \(M_0\).
\end{remark}

\subsection{A bounded family of large superlevel sets}

We record explicitly the elementary net consequence that determines
the order of the remaining choices.

\begin{lemma}[Superlevel sets from all-filter nets]
\label{lem:rare-superlevel-nets}
Consider the family of unnormalized successful base-query densities
on a prepared leaf, allowing every orientation-only filter.
Suppose a member \(d\) has total mass at least \(b_0>0\).
Uniformly over the prepared leaves and the finite formats,
there are constants \(t,\delta_0>0\) such that, for every sufficiently
small fixed \(\epsilon>0\), a bounded family of reference sets contains
a set \(W\) satisfying
\begin{equation}
 \nu(W)\ge\delta_0/2,\qquad
 \nu\{k\in W:d(k)<t/4\}\le4\epsilon/t.
 \label{eq:rare-net-superlevel}
\end{equation}
The family depends on the leaf, format, and \(\epsilon\), but not on
the orientation-only filter subsequently selected.
The constants \(t,\delta_0\) can be fixed before \(\epsilon\).
\end{lemma}

\begin{proof}
Every such density is dominated by the corresponding unfiltered
query density.  Since \(\int d\,d\nu\le1\), its superlevel set
\(\{d>T\}\) has reference measure at most \(1/T\).
The uniform small-set conclusion of \Cref{lem:histogram-control}
therefore makes the integral on this set uniformly small as \(T\)
increases.
Choose a fixed \(T\) so large that
\(\int_{\{d>T\}}d\,d\nu<b_0/4\) for every density in question
and all sufficiently large \(n\).
Put
\[
        t=b_0/4,\qquad \delta_0=b_0/(2T).
\]
The region \(t\le d\le T\) carries integral at least \(b_0/2\),
so its reference measure is at least \(\delta_0\).

Take an \(L^1(\nu)\) \(\epsilon\)-net for the entire family of
orientation-filtered densities, as supplied by
\Cref{lem:histogram-control}.  If \(f_j\) approximates \(d\), set
\(W_j=\{f_j\ge t/2\}\).
At most \(2\epsilon/t\) of the region \(d\ge t\) lies outside \(W_j\).
Consequently \(\nu(W_j)\ge\delta_0/2\) when
\(\epsilon\le\delta_0t/4\).
On \(W_j\cap\{d<t/4\}\) one has
\(\lvert f_j-d\rvert\ge t/4\), so that set has measure at most
\(4\epsilon/t\).
The sets \(W_j\) for all net representatives form the required family.
\end{proof}

\subsection{Overlap on the compatible leaf pairs}

\begin{proposition}
\label{prop:rare-overlap}
In the inverse-polynomial alternative of
\Cref{prop:status-preparation}, the overlap criterion
\eqref{eq:overlap-criterion} holds with a lower bound
\(\Omega(N^{-4})\), along the prepared sequence.
\end{proposition}

\begin{proof}
Write \(\pi_n\) for the prepared leaf distribution and
\(\sigma_\ell\) for the conditional orientation law on a leaf.
Two leaves are initially independent.
The preparation supplies a set \(\mathcal C_n\) of compatible leaf
pairs with
\[
       (\pi_n\otimes\pi_n)(\mathcal C_n)\ge c_0N^{-4}
\]
for some fixed \(c_0>0\).
For every pair in this set, choose an injecting compatible small
table whose two separate admissibility filters have orientation
masses at least \(a_0>0\).
All choices are from finite numerical table and coordinate formats;
the constants \(c_0,a_0\) do not depend on \(n\).

Choose one generic atom for each cross edge, at a common fixed tag,
with diagonal bit \(q=1\).
At each endpoint choose two distinct allowed labels.
Use role zero on the first unit and role one on the second.
Feasible summaries exist: on the first side take all ordinary-block
tester bits zero and the shared tester bit one; on the second side
take one allowed ordinary-block tester bit one, all the other such
bits zero, and the shared tester bit zero.
Prescribe the basis values at endpoint \(i\), facing endpoint \(z\),
to be
\[
             T(x_{iz},\cdot)|_{C_i}
              =(a+(u_zU_i)^*)|_{C_i}.
\]
Here the star denotes the contraction computed with the chosen table.
By \eqref{eq:rare-exact-prediction},
\[
 (p_i+a)(x_{iz})
   =T(\lambda_i,x_{iz})
   =a(\lambda_i)+(u_zU_i)^*(\lambda_i).
\]
Summing over the two endpoints \(i\) gives
\(c+u_z(\Delta_A)=0\), by the pinned compatibility.
The same calculation applies with the two units interchanged.
Thus the coupled conditions in \eqref{eq:recipe-scalars} hold,
as do its diagonal and role conditions.
Choose the binary prescriptions given by \Cref{lem:small-tables}.
These are base queries: the recorded old marks and prescribed basis
values already determine the \(p_i\)-values, so no additional
\(p_i\)-filter is imposed.
By \Cref{thm:gradient-realization}, equal accepted keys therefore
satisfy the hypotheses of \Cref{prop:collision}.

At all but an exponentially small proportion of orientations,
\Cref{lem:binary-mixing,lem:unary-positivity} give a fixed positive
lower bound for the total successful mass of every feasible numerical
base request.  There are finitely many requests.
Markov's inequality over conditional leaf probabilities shows that,
outside an exponentially small set of second-side leaves, every
admissibility filter of mass at least \(a_0\) retains total successful
query mass at least a fixed \(b_0>0\).
For example, first make the conditional fraction of exceptional
orientations at most \(a_0/2\); the fixed success lower bound times
\(a_0/2\) then gives \(b_0\).
Removing these leaves loses \(o(N^{-4})\) leaf-pair mass even before
compatibility is imposed.

Apply \Cref{lem:rare-superlevel-nets} on each remaining B-leaf
and each numerical base format.
Fix its uniform constants \(t,\delta_0\), and put
\[
 \delta=\delta_0/2,\qquad
 c_{\mathrm{hit}}=c_{\mathrm{hit}}(\delta),\qquad
 A_0=(a_0/2)c_{\mathrm{hit}}>0.
\]
Use uniform per-leaf small-set control to choose \(\eta>0\) such that
every reference set of measure at most \(\eta\) has A-query mass
less than \(A_0/2\), for all sufficiently large \(n\).
This applies also to a filtered query, which is dominated by the
unfiltered one.  Now fix a net accuracy satisfying
\[
          \epsilon\le
          \min\{\delta_0t/4,\eta t/4\}.
\]
For each B-leaf and numerical base format, take the family of sets
\(W_{\ell_B,r,j}\) supplied at this accuracy, and let \(L\) be a
uniform bound on its cardinality.  The constants through \(L\) are
independent of \(n\), although the sets themselves may vary with \(n\)
and the B-leaf.  Each family was formed over all orientation-only filters.
For a compatible pair, its filtered B-density \(d_B\) has mass at
least \(b_0\), so choose a member \(W\) satisfying
\eqref{eq:rare-net-superlevel} for that density.
Although the table and filter may depend on \(\ell_A\), that
dependence only selects a member of the already fixed B-leaf family.

We now quantify the first-side exceptional pairs.
Let \(C\) bound the number of numerical request and shape choices.
Discard family members of reference measure below \(\delta\);
the selected member is not discarded.
Let \(e_n\) denote the supremum bad probability in
\eqref{eq:rare-superpolynomial-tail}, over the remaining sets and
all relevant requests.
For a fixed B-leaf every such set is deterministic for an independent
A-orientation.  A union bound and then independence of the original
leaves give
\begin{equation}
 \E_{\ell_A,\ell_B}
 \PP_{\omega_A\mid\ell_A}
 \bigl(\text{some relevant }W_{\ell_B,r,j}
       \text{ fails its feasible hit bound}\bigr)
 \le CL e_n .
 \label{eq:rare-independent-net-error}
\end{equation}
Hence the leaf pairs on which the inner probability exceeds \(a_0/2\)
have mass at most \(2CL e_n/a_0=o(N^{-4})\), by
\Cref{lem:rare-large-set-hit}.
On any other compatible pair, at least \(a_0/2\) of the A-orientations
pass the admissibility filter and all relevant hit bounds.  For its
selected \(W\) and unnormalized density \(d_A\), this gives
\begin{equation}
       \int_W d_A\,d\nu\ge A_0.
 \label{eq:rare-a-hit}
\end{equation}

Together with the exceptional B-leaves, the discarded leaf pairs
consume \(o(N^{-4})\) mass.

On a remaining compatible pair let
\(H=\{k\in W:d_B(k)<t/4\}\).
Then \(\nu(H)\le\eta\), and \eqref{eq:rare-a-hit} yields
\[
 \int d_A d_B\,d\nu
   \ge \frac t4\left(\int_Wd_A\,d\nu-\int_Hd_A\,d\nu\right)
   \ge \frac{tA_0}{8}.
\]
There remain at least \((c_0/2)N^{-4}\) compatible pairs for large \(n\).
Consequently
\begin{equation}
       \E_{\ell_A,\ell_B}\int d_A d_B\,d\nu
          \ge \frac{c_0tA_0}{16}N^{-4}.
 \label{eq:rare-overlap-bound}
\end{equation}
Tables and numerical requests may be chosen separately for each
leaf pair, as allowed in \Cref{prop:collision}; alternatively one may
sum over their finitely many options.
This proves the overlap criterion.
\end{proof}

\subsection{Completion}

\begin{proof}[Proof of \Cref{thm:raw-supersaturation}]
Fix the construction constants and, for every sufficiently large \(n\),
fix mixers satisfying \Cref{lem:generic-mixers}, before choosing a unit law.
Suppose there is no common sufficiently-large-\(n\) bound as asserted.
Then along an increasing sequence of \(n\)'s there are admissible
unit laws \(\sigma_n\) whose probability of four cross holes is
smaller than \(2^{-100gN}\).
Apply \Cref{lem:peeling,prop:status-preparation}, taking subsequences
where needed.  All retained laws have relative mass \(2^{-o(N)}\),
and their final leaves satisfy the required image bounds.

Every alternative with positive limiting mass of injecting compatible
tables gives the overlap criterion by \Cref{prop:positive-overlap}.
The remaining inverse-polynomial alternative gives it by
\Cref{prop:rare-overlap}.
In both cases \Cref{prop:collision}, including the cost of returning
to the original laws, gives four-hole probability at least
\[
      2^{-(k_{\max}+.03)N-o(N)}.
 \]
But
\[
 100g-k_{\max}-.03
   =100g-56(g-1)-.03
   =44g+55.97>0.
\]
The displayed lower bound is therefore larger than \(2^{-100gN}\)
for large \(n\), a contradiction.
Since a failure of uniformity over admissible laws produced precisely
such a sequence, the resulting sufficiently-large-\(n\) statement
is uniform over those laws.
\end{proof}

\begin{proof}[Proof of \Cref{thm:main}]
Apply \Cref{prop:raw-to-graph} to \Cref{thm:raw-supersaturation}.
For every sufficiently large admissible \(n\), there is a graph on
\(m=2^{C_0gN}\) sampled positions with independence number at most two
and with \(\cm(G)<m/100\).
Since all constants, including \(M_0\), are fixed and
\(N=M_0n\) tends to infinity, these orders are arbitrarily large.
\end{proof}
